% from ega0-21-fr.tex
\oldpage[0\textsubscript{IV-1}]{177}
$n=\operatorname{rg}_K(\Omega^1_{A/k_0}\otimes_AK)$. En vertu de (21.9.1),
$\widehat\Omega^1_{A/k_0}$ est un $A$-module de type fini, et en outre ce
$A$-module est plat en vertu de ($\mathbf 0_{\mathrm{III}}$, 10.2.1)~; il est
donc libre ($\mathbf 0_{\mathrm{III}}$, 10.1.3) de rang $n$, d'où la première
assertion.

Comme par hypothèse $\Omega^1_{k_0/k}$ est de rang fini sur $k_0$, le produit
tensoriel complété
$\Omega^1_{k_0/k}\widehat\otimes_{k_0}A$ s'identifie à
$\Omega^1_{k_0/k}\otimes_{k_0}A$~; comme $A$ est une $k_0$-algèbre
formellement lisse, il résulte de (20.7.18) que l'homomorphisme
\[
  \widehat u_{A/k_0/k}:
  \Omega^1_{k_0/k}\otimes_{k_0}A
  \longrightarrow\widehat\Omega^1_{A/k}
\]
admet un inverse à gauche qui est un $A$-homomorphisme continu, ce qui
implique en particulier que son image est fermée dans
$\widehat\Omega^1_{A/k}$~; la suite exacte
\[
  0\longrightarrow \Omega^1_{k_0/k}\otimes_{k_0}A
  \xrightarrow{\widehat u_{A/k_0/k}}\widehat\Omega^1_{A/k}
  \longrightarrow\widehat\Omega^1_{A/k_0}\longrightarrow0
\]
est donc exacte et scindée en vertu de (20.7.17) et de ce qui précède~; ce
qui achève de prouver la proposition.
\end{proof}

\begin{corollary}[21.9.3]
\label{0.21.9.3-fr}
Soient $k$ un corps, $A$ l'anneau de séries formelles
$k[[T_1,\ldots,T_r]]$, muni de sa structure usuelle de $k$-algèbre. Alors
$\widehat\Omega^1_{A/k}$ est un $A$-module libre de rang $r=\dim(A)$.
\end{corollary}

\begin{proof}
Il suffit de noter que $A$ est une $k$-algèbre formellement lisse (19.3.4),
et d'appliquer (21.9.2) avec $k_0=k=K$.
\end{proof}

\begin{lemma}[21.9.4]
\label{0.21.9.4-fr}
Soient $k$ un corps de caractéristique $p>0$, $A$ une $k$-algèbre qui est un
anneau local noethérien complet, $B$ une sous-$k$-algèbre de $A$ isomorphe à
une $k$-algèbre topologique de la forme $k[[T_1,\ldots,T_r]]$ et telle que
$A$ soit une $B$-algèbre de type fini. Si $B_1$ est la sous-algèbre
$k[[T_1^p,\ldots,T_r^p]]$ de $B$, $\widehat\Omega^1_{A/k}$ s'identifie
canoniquement à $\Omega^1_{A/B_1}$.
\end{lemma}

\begin{proof}
Toute dérivation continue de $B_1$ dans un $A$-module topologique $P$, qui
est restriction d'une $k$-dérivation continue de $A$ dans $P$, est \emph{nulle},
car elle l'est dans le sous-anneau $k[T_1^p,\ldots,T_r^p]$ des polynômes, et
ce dernier est dense dans $B_1$. La suite exacte (20.3.7.1) montre donc que
l'homomorphisme canonique
\[
  \operatorname{D\acute{e}r.\,cont}_{B_1}(A,P)
  \longrightarrow \operatorname{D\acute{e}r.\,cont}_k(A,P)
\]
est bijectif~; la conclusion résulte de (20.7.14.4), compte tenu de ce que
$\Omega^1_{A/B_1}$ est un $A$-module de type fini (20.4.7), donc séparé et
complet puisque $A$ est un anneau local noethérien complet.
\end{proof}

\begin{proposition}[21.9.5]
\label{0.21.9.5-fr}
Soient $A$ un anneau local noethérien complet intègre, $A_0$ un sous-anneau
de $A$ isomorphe à un anneau de séries formelles $k[[T_1,\ldots,T_r]]$ sur
un corps $k$, tel que $A$ soit une $A_0$-algèbre finie et que le corps des
fractions $E$ de $A$ soit une extension séparable du corps des fractions
$L_0$ de $A_0$. Alors on a
\begin{equation}
\tag{21.9.5.1}
\label{0.21.9.5.1-fr}
  \operatorname{rg}_A\widehat\Omega^1_{A/k}
  =\operatorname{rg}_E(\widehat\Omega^1_{A/k}\otimes_AE)=\dim(A).
\end{equation}
\end{proposition}

\begin{proof}
Notons que si $\mathfrak m_0$ est l'idéal maximal de $A_0$, la topologie de
$A$ est la topologie $\mathfrak m_0$-adique puisque $A$ est une
$A_0$-algèbre finie (Bourbaki, \emph{Alg. comm.}, chap.~IV, \S~2,
n\textsuperscript{o}~5, cor.~3 de la prop.~9) et induit sur $A_0$ la topologie
$\mathfrak m_0$-adique (Bourbaki, \emph{Alg. comm.}, chap.~III, \S~3,
n\textsuperscript{o}~4, th.~3). On sait (21.9.1) que
$\widehat\Omega^1_{A/k}$ est un $A$-module de type fini et,

